% 由 tools/paper-results.py 生成，请勿手改。
\begin{tabularx}{\textwidth}{@{}llLrrrrrrr@{}}
\toprule
\bt{Class}{类别} & \bt{Change}{数据变化} & \bt{What the update does}{更新内容} & \bh{\swatch{mNoShare}\\No\\sharing}{不共享} & \bh{\swatch{mTraj}\\Trajectory\\retrieval}{轨迹检索} & \bh{\swatch{mUnguarded}\\Un-\\guarded}{无守护} & \bh{\swatch{mSchema}\\Schema-\\only}{只看结构} & \bh{\swatch{mRevoke}\\Revoke-\\on-write}{写入即撤销} & \bh{\swatch{mDef}\\Definition-\\level}{定义级} & \bh{\swatch{mCond}\\\system\\(condition)}{条件级}\\
\midrule
\bt{None}{无变化} & \bt{Held-out}{留出} & \bt{No update; new parameters and question types}{无更新；新参数与新题型} & \cellcolor{heat!30}49 & \cellcolor{heat!56}100 & \cellcolor{heat!56}100 & \cellcolor{heat!56}100 & \cellcolor{heat!54}96 & \cellcolor{heat!52}91 & \cellcolor{heat!56}100\\
\addlinespace[2pt]
\bt{Benign}{正常} & \bt{Append}{正常追加} & \bt{Copy one month of sales under new tickets}{复制一个月销售并换新小票号} & \cellcolor{heat!25}38 & \cellcolor{heat!56}100 & \cellcolor{heat!54}96 & \cellcolor{heat!56}100 & \cellcolor{heat!47}81 & \cellcolor{heat!47}81 & \cellcolor{heat!54}96\\
 & \bt{Late backfill}{迟到回填} & \bt{Add returns with new keys}{补写新键的退货} & \cellcolor{heat!27}41 & \cellcolor{heat!56}100 & \cellcolor{heat!56}100 & \cellcolor{heat!56}100 & \cellcolor{heat!31}50 & \cellcolor{heat!39}67 & \cellcolor{heat!56}100\\
 & \bt{In-place fix}{原地更正} & \bt{Update net paid of some sales}{原地修改部分销售的净支付额} & \cellcolor{heat!21}30 & \cellcolor{heat!56}100 & \cellcolor{heat!56}100 & \cellcolor{heat!56}100 & \cellcolor{heat!52}93 & \cellcolor{heat!49}85 & \cellcolor{heat!56}100\\
 & \bt{New column}{新增无关列} & \bt{Add an empty column to returns}{退货表新增空列} & \cellcolor{heat!39}67 & \cellcolor{heat!56}100 & \cellcolor{heat!56}100 & \cellcolor{heat!53}94 & \cellcolor{heat!50}89 & \cellcolor{heat!45}78 & \cellcolor{heat!56}100\\
\addlinespace[2pt]
\bt{Grain}{粒度} & \bt{Status history}{状态流水} & \bt{Add an application-state row per return}{每笔退货追加申请状态行} & \cellcolor{heat!25}39 & \cellcolor{heat!31}50 & \cellcolor{heat!31}50$^\dagger$ & \cellcolor{heat!31}50$^\dagger$ & \cellcolor{heat!45}78 & \cellcolor{heat!45}78 & \cellcolor{heat!56}100\\
 & \bt{Versioning}{版本化更正} & \bt{Keep old sales versions as non-current rows}{旧版本销售保留为非当前行} & \cellcolor{heat!17}22 & \cellcolor{heat!45}78 & \cellcolor{heat!45}78$^\dagger$ & \cellcolor{heat!43}74$^\dagger$ & \cellcolor{heat!49}85 & \cellcolor{heat!49}85 & \cellcolor{heat!54}96\\
 & \bt{Duplicate load}{重复装载} & \bt{Reload four months of identical sales rows}{四个月销售整批重复装载} & \cellcolor{heat!17}22 & \cellcolor{heat!25}37 & \cellcolor{heat!29}46$^\dagger$ & \cellcolor{heat!28}44$^\dagger$ & \cellcolor{heat!17}22 & \cellcolor{heat!30}48 & \cellcolor{heat!36}59\\
\addlinespace[2pt]
\bt{Fan-out}{连接放大} & \bt{Dim.\ history}{维表拉链} & \bt{Add old-version rows to the item dimension}{商品维表增加旧版本行} & \cellcolor{heat!6}0 & \cellcolor{heat!56}100 & \cellcolor{heat!56}100$^\dagger$ & \cellcolor{heat!56}100$^\dagger$ & \cellcolor{heat!56}100 & \cellcolor{heat!56}100 & \cellcolor{heat!56}100\\
\addlinespace[2pt]
\bt{Coverage}{覆盖} & \bt{Key format}{日期键格式} & \bt{Append sales whose date keys use yyyymmdd}{追加日期键为 yyyymmdd 的销售} & \cellcolor{heat!12}11 & \cellcolor{heat!17}22 & \cellcolor{heat!17}22$^\dagger$ & \cellcolor{heat!17}22$^\dagger$ & \cellcolor{heat!15}19 & \cellcolor{heat!12}11 & \cellcolor{heat!12}11\\
\addlinespace[2pt]
\bt{Not modeled}{未建模} & \bt{Unit change}{金额单位} & \bt{Multiply sales amounts by 100 from 2002-07}{2002-07 起销售金额乘以 100} & \cellcolor{heat!17}22 & \cellcolor{heat!30}48 & \cellcolor{heat!27}42$^\dagger$ & \cellcolor{heat!31}50$^\dagger$ & \cellcolor{heat!28}44 & \cellcolor{heat!25}37$^\dagger$ & \cellcolor{heat!28}44$^\dagger$\\
\addlinespace[2pt]
\bt{Ambiguous fix}{修复有歧义} & \bt{Backup copy}{备份副本} & \bt{Append a full backup copy that keeps old amounts}{追加保留旧金额的整份备份副本} & \cellcolor{heat!28}44 & \cellcolor{heat!42}72 & \cellcolor{heat!39}67$^\dagger$ & \cellcolor{heat!37}61$^\dagger$ & \cellcolor{heat!28}44 & \cellcolor{heat!37}61 & \cellcolor{heat!37}61\\
\midrule
\multicolumn{3}{@{}l}{\bt{Mean over the 12 settings}{12 种情形平均}} & \textbf{32} & \textbf{76} & \textbf{75} & \textbf{75} & \textbf{67} & \textbf{69} & \textbf{81}\\
\bottomrule
\end{tabularx}
